\documentclass[11pt]{amsart}

\usepackage[T1]{fontenc}
\usepackage{lmodern}
\usepackage{amsmath,amssymb,amsthm,mathtools}
\usepackage{booktabs,tabularx,array,microtype}
\usepackage{xcolor}
\usepackage[colorlinks=true,linkcolor=blue!55!black,citecolor=blue!55!black,urlcolor=blue!55!black]{hyperref}

\allowdisplaybreaks
\setlength{\emergencystretch}{1.5em}

\newtheorem{theorem}{Theorem}[section]
\newtheorem{lemma}[theorem]{Lemma}
\newtheorem{proposition}[theorem]{Proposition}
\newtheorem{corollary}[theorem]{Corollary}
\theoremstyle{definition}
\newtheorem{definition}[theorem]{Definition}
\newtheorem{remark}[theorem]{Remark}

\newcommand{\Q}{\mathbb{Q}}
\newcommand{\Z}{\mathbb{Z}}
\newcommand{\phig}{\frac{1+\sqrt5}{2}}

\title[Non-Parry powers of the quartic Perron root]{Every Power of the Perron Root of
\texorpdfstring{$X^4-X-1$}{X4-X-1} Is Non-Parry:\ An Exact Certificate}
\author{Zijian Zeng}
\address{Institute of Computer Science and Digital Innovation, UCSI University, Kuala Lumpur 56000, Malaysia}
\email{zijianzeng@foxmail.com}
\email{1002266693@ucsiuniversity.edu.my}
\date{August 30, 2026}
\subjclass[2020]{Primary 11A63; Secondary 11R06, 37B10}
\keywords{beta-expansion, Parry number, Perron number, algebraic certificate}
\hypersetup{
  pdftitle={Every Power of the Perron Root of X4-X-1 Is Non-Parry: An Exact Certificate},
  pdfauthor={Zijian Zeng}
}

\begin{document}

\begin{abstract}
Let $\beta>1$ be the real root of $X^4-X-1$.  We prove that
$\beta^k$ is non-Parry for every integer $k\geq 1$.  This settles the fixed
$d=4$ case of a conjecture of Hare and Hichri, who recorded this case as part
of broader computational evidence.  A nonreal conjugate $\gamma$ gives a
uniform obstruction for $k\geq 8$ through Solomyak's golden-ratio bound.  The
seven remaining powers are treated by exact greedy prefixes whose embedded
remainders escape every eventually periodic tail bound.  All floor decisions
and strict inequalities are reduced to rational interval calculations.  A
serialized certificate is checked independently by a standard-library
rational verifier and by a symbolic verifier using Sturm root counts.
\end{abstract}

\maketitle

\section{Introduction}

For a real base $q>1$, let
\[
T_q(x)=qx-\lfloor qx\rfloor \qquad (0\leq x\leq 1)
\]
and define the greedy digits of exact $1$ by
\[
d_n(q)=\left\lfloor qT_q^{n-1}(1)\right\rfloor,\qquad n\geq 1.
\]
The base $q$ is a \emph{Parry number} when the word
$d_1(q)d_2(q)\cdots$ is eventually periodic; a finite expansion is understood
to have an all-zero tail.  This is equivalent to finiteness of the orbit of
$1$ under $T_q$ and is the standard convention originating with
Parry~\cite{Parry1960}.

Let $\beta_d>1$ be the Perron root of $X^d-X-1$.  Akiyama proved that
$\beta_d$ is non-Parry for every $d\geq 4$~\cite[Theorem~1.1]{Akiyama2016}.
Hare and Hichri subsequently conjectured that every positive power
$\beta_d^k$ is non-Parry for all $d\geq 4$~\cite[Conjecture~3.3]{HareHichri2026}.
Their Table~3.3 labels the cases $4\leq d\leq 10$ as non-Parry for all positive
powers, while the surrounding text describes the table as numerical support
and reports computations through $d=100$~\cite[Table~3.3 and
Remark~3.4]{HareHichri2026}.  The exact finite witnesses behind the $d=4$ row
are not supplied there.

The purpose of this note is to give a fully exact and independently
reproducible proof for that fixed row.

\begin{theorem}\label{thm:main}
Let $\beta>1$ satisfy $\beta^4=\beta+1$.  Then $\beta^k$ is non-Parry for
every integer $k\geq 1$.
\end{theorem}

The proof has two parts.  A nonreal conjugate $\gamma$ satisfies
$|\gamma|>1$ and $|\gamma|^8>\phig$.  Solomyak's strict bound on conjugates of
Parry numbers~\cite{Solomyak1994} therefore covers every $k\geq 8$.  For
$1\leq k\leq 7$, we give a short exact greedy prefix and prove that its image
under $\beta\mapsto\gamma$ is already too large to be the tail of an
eventually periodic expansion.  This escape principle is related to
Akiyama's criterion~\cite[Lemma~2.1]{Akiyama2016}; we rederive it for the
exact-$1$ convention used here.

\section{Algebraic root data}

Put
\[
f(X)=X^4-X-1,
\qquad
h(S)=S^6+S^4-S^3-S^2-1.
\]
We use the rational intervals
\begin{align}
B_-&=\frac{1220744084}{10^9}, &
B_+&=\frac{1220744085}{10^9},\label{eq:Binterval}\\
S_-&=\frac{1130685445}{10^9}, &
S_+&=\frac{1130685446}{10^9}.\label{eq:Sinterval}
\end{align}

\begin{lemma}\label{lem:roots}
The polynomial $f$ is irreducible over $\Q$.  It has a unique positive root
$\beta\in(B_-,B_+)$.  If $s$ is the unique positive root of $h$ in
$(S_-,S_+)$ and
\[
u=-\frac{s}{s^2+1},\qquad p=-\frac1s,
\]
then
\begin{equation}\label{eq:factor}
f(X)=(X^2-uX+s)(X^2+uX+p).
\end{equation}
The first quadratic has conjugate nonreal roots $\gamma,\overline\gamma$ with
$|\gamma|^2=s$; the second has roots $\beta$ and a real $\alpha$ with
$|\alpha|<1$.  Moreover,
\[
1<|\gamma|<\beta.
\]
\end{lemma}

\begin{proof}
Modulo $2$, the polynomial $f$ becomes $X^4+X+1$.  It has no root in
$\mathbb F_2$ and has remainder $1$ upon division by the only monic
irreducible quadratic $X^2+X+1$.  Hence it is irreducible.

Direct rational substitution gives
$f(B_-)<0<f(B_+)$ and $h(S_-)<0<h(S_+)$.  Each polynomial has one sign change,
so Descartes' rule gives uniqueness of its positive root.  Since
\[
(s^2-1)(s^2+1)^2-s^3=h(s)=0,
\]
we have $p+s-u^2=0$, $u(s-p)=-1$, and $sp=-1$.  Multiplication now proves
\eqref{eq:factor}.

The discriminant of the first factor is negative because
$u^2=s-1/s<4s$.  Its roots have product $s$, hence modulus $\sqrt{s}$.  The
second factor has negative product and therefore one positive and one negative
root; its positive root is the unique positive root $\beta$ of $f$.  The
intervals give $S_+<B_-^2$ and $S_-B_->1$.  Thus
$|\gamma|^2=s<\beta^2$, while
$|\alpha|=1/(s\beta)<1$.
\end{proof}

\begin{lemma}\label{lem:cutoff}
For the conjugate $\gamma$ in Lemma~\ref{lem:roots},
\[
|\gamma|^8>\phig.
\]
\end{lemma}

\begin{proof}
Let $c=16181/10000$.  Exact rational arithmetic gives
\[
s^4>S_-^4>c.
\]
Also $c>1/2$ and $(2c-1)^2>5$, so $c>\phig$.
\end{proof}

\section{An escape criterion for exact 1}

\begin{lemma}[embedded-tail escape]\label{lem:escape}
Let $q>1$ be algebraic, and let $d_1d_2\cdots$ be the greedy expansion of
exact $1$ in base $q$.  Define
\[
r_0=1,\qquad r_n=qr_{n-1}-d_n.
\]
Suppose an embedding of a number field containing $q$ sends $q$ to
$\eta$, where $|\eta|>1$, and let $r_n'$ be the image of $r_n$.  If
\begin{equation}\label{eq:escape}
|r_n'|>\frac{\lfloor q\rfloor}{|\eta|-1}
\end{equation}
for some $n$, then $q$ is non-Parry.
\end{lemma}

\begin{proof}
Assume that the digit word is eventually periodic.  The identity
\[
\sum_{i\geq1}d_iq^{-i}=1
\]
is then a rational function identity in $q^{-1}$, so applying the embedding is
legitimate.  Since $|\eta|>1$, the corresponding
eventually periodic series converges at $\eta^{-1}$ and represents that same
rational function.  Consequently,
\[
r_n'=\eta^n\left(1-\sum_{i=1}^n d_i\eta^{-i}\right)
     =\sum_{j\geq1}d_{n+j}\eta^{-j}.
\]
Every greedy digit lies between $0$ and $\lfloor q\rfloor$, so
\[
|r_n'|\leq \sum_{j\geq1}\lfloor q\rfloor |\eta|^{-j}
=\frac{\lfloor q\rfloor}{|\eta|-1},
\]
contrary to \eqref{eq:escape}.  A finite expansion is included by its zero
tail.
\end{proof}

\section{The seven finite certificates}

All states are computed in the free rank-four ring
\[
\Z[X]/(X^4-X-1),
\]
using $X^4=X+1$.  For fixed $k$, start with $R_0(X)=1$ and use the digits in
Table~\ref{tab:witnesses} to form
\begin{equation}\label{eq:recurrence}
R_j(X)\equiv X^kR_{j-1}(X)-d_j \pmod{f(X)},
\qquad \deg R_j\leq 3.
\end{equation}

\begin{table}[t]
\centering
\caption{Exact greedy-prefix escape data.  The last state is
$R_n(\beta)=P_k(\beta)$.}\label{tab:witnesses}
\small
\setlength{\tabcolsep}{7pt}
\begin{tabular}{@{}ccclc@{}}
\toprule
$k$ & $\lfloor\beta^k\rfloor$ & $n$ & $d_1\cdots d_n$ & $L_k$\\
\midrule
1 & 1 & 35 & \texttt{10000000100000000000010000000010000} & $16$\\
2 & 1 & 10 & \texttt{1010000001} & $8$\\
3 & 1 & 7  & \texttt{1101100} & $5$\\
4 & 2 & 7  & \texttt{2010002} & $17/2$\\
5 & 2 & 3  & \texttt{212} & $6$\\
6 & 3 & 3  & \texttt{310} & $15/2$\\
7 & 4 & 4  & \texttt{4002} & $10$\\
\bottomrule
\end{tabular}
\end{table}

The corresponding final states are
\begin{align*}
P_1(X)&=7+6X-7X^2-2X^3,&
P_2(X)&=-4-3X+3X^2+2X^3,\\
P_3(X)&=1-X-2X^2+2X^3,&
P_4(X)&=1+5X-X^2-3X^3,\\
P_5(X)&=-1+3X+X^2-2X^3,&
P_6(X)&=3+X-4X^2+X^3,\\
P_7(X)&=-5+5X^2-X^3.&&
\end{align*}

We now describe a finite exact check that certifies both greediness and escape.
For $P(X)=\sum_{i=0}^3a_iX^i$, define the outward rational enclosure
\begin{align*}
\underline P&=a_0+
 \sum_{\substack{1\leq i\leq3\\a_i\geq0}}a_iB_-^i+
 \sum_{\substack{1\leq i\leq3\\a_i<0}}a_iB_+^i,\\
\overline P&=a_0+
 \sum_{\substack{1\leq i\leq3\\a_i\geq0}}a_iB_+^i+
 \sum_{\substack{1\leq i\leq3\\a_i<0}}a_iB_-^i.
\end{align*}
Then $P(\beta)\in[\underline P,\overline P]$.  Applying this enclosure to
every state generated by \eqref{eq:recurrence} gives
\begin{equation}\label{eq:statebounds}
0<\underline {R_j}\leq R_j(\beta)\leq\overline {R_j}<1
\end{equation}
in all seven rows.  Thus the displayed strings are exact greedy prefixes.
The same calculation gives the listed values of $\lfloor\beta^k\rfloor$.

It remains to certify the embedded norms without approximating $\gamma$.
For a final state $P_k(X)=a_0+a_1X+a_2X^2+a_3X^3$, reduce with
$\gamma^2=u\gamma-s$ to write
\[
P_k(\gamma)=A_k+B_k\gamma,
\]
where
\begin{align*}
A_k&=a_0-a_2s-a_3us,\\
B_k&=a_1+a_2u+a_3(u^2-s).
\end{align*}
It follows that
\begin{equation}\label{eq:norm}
|P_k(\gamma)|^2=A_k^2+A_kB_ku+B_k^2s.
\end{equation}
After subtracting $L_k^2$, clearing the positive denominator, and reducing
the numerator modulo $h(s)$, one obtains the polynomials $G_k$ listed in
Appendix~\ref{app:gaps}.  Direct outward rational evaluation on
$[S_-,S_+]$ proves
\begin{equation}\label{eq:normlower}
G_k(s)>0,\qquad |P_k(\gamma)|>L_k
\quad (1\leq k\leq7).
\end{equation}

Finally, put
\[
\rho=\frac{1063336938}{10^9}.
\]
The exact inequalities $\rho^2<S_-$ and
\begin{equation}\label{eq:products}
L_k(\rho^k-1)>\lfloor\beta^k\rfloor
\qquad(1\leq k\leq7)
\end{equation}
are rational comparisons.  Since $|\gamma|>\rho$, equations
\eqref{eq:normlower} and \eqref{eq:products} give
\[
|P_k(\gamma)|>
\frac{\lfloor\beta^k\rfloor}{|\gamma|^k-1}.
\]
Lemma~\ref{lem:escape}, applied with $q=\beta^k$ and $\eta=\gamma^k$,
proves that $\beta^k$ is non-Parry for $1\leq k\leq7$.

\section{Completion of the proof}

\begin{proof}[Proof of Theorem~\ref{thm:main}]
The preceding section covers $1\leq k\leq7$.  For $k\geq8$, restriction of
the embedding $\Q(\beta)\hookrightarrow\mathbb C$, $\beta\mapsto\gamma$, to
$\Q(\beta^k)$ maps $\beta^k$ to $\gamma^k$.  These images are distinct because
$|\gamma|<\beta$, so $\gamma^k$ is a non-dominant conjugate of $\beta^k$.
Lemma~\ref{lem:cutoff} gives
\[
|\gamma^k|\geq|\gamma|^8>\phig.
\]
Solomyak proved that every non-dominant Galois conjugate of a Parry number has
modulus strictly less than $\phig$~\cite{Solomyak1994}.  Hence $\beta^k$ is
non-Parry for every $k\geq8$ as well.
\end{proof}

\begin{remark}
The field $\Q(\beta^k)$ is allowed to be smaller than $\Q(\beta)$ in the last
argument.  Restricting an embedding still produces a conjugate of $\beta^k$;
no degree-preservation assertion is needed.
\end{remark}

\section{Exact verification and limitations}

The ancillary certificate contains only integers and rationals: the intervals
\eqref{eq:Binterval}--\eqref{eq:Sinterval}, the seven words in
Table~\ref{tab:witnesses}, the final states, the $L_k$, and the $G_k$.  It is
verified in two independent ways.

\begin{enumerate}
\item \texttt{certify\_exact.py} uses the Python standard library only.  It
checks irreducibility, isolating signs, the factorization, every state bound
in \eqref{eq:statebounds}, equation \eqref{eq:norm}, the seven strict products,
and the exponent-eight cutoff by outward rational interval arithmetic.
\item \texttt{verify\_symbolic.py} independently recomputes quotient-ring
remainders and uses exact Sturm root counts for every sign.  It derives the
norm numerators symbolically and checks their remainders modulo $h$.
\end{enumerate}

From the ancillary directory, the checks are run by
\begin{verbatim}
python3 certify_exact.py nonparry_d4.json
python3 verify_symbolic.py nonparry_d4.json
\end{verbatim}
The second command was tested with Python 3.10.12 and SymPy 1.14.0; the first
has no third-party dependency.  Neither command reads floating-point discovery
output.

The theorem resolves only the fixed $d=4$ case.  It does not settle the
Hare--Hichri conjecture for all $d\geq4$.  Searches of the arXiv, DOI records,
and exact certificate fingerprints through August 30, 2026 found the prior
computational report in~\cite{HareHichri2026}, but no published exact witness
data for this row.  This is a bounded literature statement, not a claim about
unindexed or unpublished work.

\section*{Data, code, and computational disclosure}

The LaTeX source, exact JSON certificate, and both verification programs are
included with the submission archive.  Floating-point computation was used
only to locate candidate prefixes; it is not imported by either checker and is
not used in the proof.  Generative AI assistance was used during exploration,
code drafting, adversarial checking, and manuscript preparation.  The stated
theorem is supported by the human-readable proof and the independently
recomputed exact certificate.  No proof assistant was used.

\appendix
\section{Norm-gap polynomials}\label{app:gaps}

With the notation of \eqref{eq:norm}, the numerator of
$|P_k(\gamma)|^2-L_k^2$, after clearing a positive denominator and reducing
modulo $h(s)$, is
\begin{align*}
G_1(S)&=281S^5-193S^4+385S^3-581S^2+220S-264,\\
G_2(S)&=71S^5-50S^4+97S^3-146S^2+56S-65,\\
G_3(S)&=29S^5-17S^4+42S^3-54S^2+23S-24,\\
G_4(S)&=320S^5-210S^4+447S^3-644S^2+256S-302,\\
G_5(S)&=42S^5-25S^4+60S^3-79S^2+35S-37,\\
G_6(S)&=272S^5-142S^4+395S^3-472S^2+212S-210,\\
G_7(S)&=132S^5-54S^4+200S^3-194S^2+107S-84.
\end{align*}
Exact interval evaluation on $[S_-,S_+]$ gives positive lower bounds for all
seven values.  These identities and signs are recomputed, rather than trusted
as input, by the symbolic verifier.

\clearpage
\begingroup
\sloppy
\bibliographystyle{amsplain}
\bibliography{references}
\endgroup

\end{document}
